\section{Lecture 23}

Recall that the ridge regression estimator is
\[
    \hat{\beta}(\lambda)
    = (X^{\top} X + \lambda I_p)^{-1} X^{\top} Y.
\]

\subsection{Properties of the Ridge Regression Estimator}

Assume $\varepsilon \sim N_n(0, \sigma^2 I)$.

\begin{result}
\[
    \mathbb{E}[\hat{\beta}(\lambda)]
    = (X^{\top} X + \lambda I)^{-1} X^{\top} \mathbb{E}[Y]
    = (X^{\top} X + \lambda I)^{-1} X^{\top} X \beta
    = \beta - \lambda (X^{\top} X + \lambda I)^{-1} \beta.
\]
Thus, $\hat{\beta}(\lambda)$ is \textbf{biased}.
\end{result}

\begin{result}
\[
    \Var(\hat{\beta}(\lambda))
    = \sigma^2 (X^{\top} X + \lambda I)^{-1} X^{\top}
        X (X^{\top} X + \lambda I)^{-1}.
\]
\end{result}

\begin{result}
It can be shown that
\[
    \Var(\hat{\beta}_{LS}) - \Var(\hat{\beta}(\lambda))
\]
is a nonnegative definite matrix.  
Hence ridge regression has \textbf{smaller variance} than least squares.
\end{result}

\begin{result}
$\hat{\beta}(\lambda)$ is multivariate normal with mean and variance as above.
\end{result}

\subsection{Mean-Square Error for a Multivariate Estimator}

Suppose $\hat{\theta}$ estimates a $p$-dimensional parameter $\theta$.
\[
    \mathrm{MSE}(\hat{\theta})
    = \mathbb{E}[(\hat{\theta} - \theta)^{\top}(\hat{\theta} - \theta)].
\]

Expanding,
\[
    \mathrm{MSE}(\hat{\theta})
    = \sum_{i=1}^p \Var(\hat{\theta}_i)
        + \|\mathrm{Bias}(\hat{\theta})\|^2
    = \operatorname{tr}(\Var(\hat{\theta})) + \|\mathrm{Bias}(\hat{\theta})\|^2.
\]

\subsection{Degrees of Freedom in Ridge Regression}

Define fitted values
\[
    \hat{Y}(\lambda)
    = X \hat{\beta}(\lambda)
    = X (X^{\top} X + \lambda I)^{-1} X^{\top} Y
    = H_{\lambda} Y,
\]
where
\[
    H_{\lambda}
    = X (X^{\top} X + \lambda I)^{-1} X^{\top}.
\]

For $\lambda = 0$, $H_0 = P_X$, the projection matrix onto $\operatorname{col}(X)$.

\begin{definition}
The \textbf{effective degrees of freedom} of ridge regression is
\[
    df(\lambda) = \operatorname{tr}(H_{\lambda}).
\]
\end{definition}

\begin{result}
If $d_1, \dots, d_p$ are the singular values of $X$, then
\[
    df(\lambda)
    = \sum_{j=1}^p \frac{d_j^2}{d_j^2 + \lambda}.
\]
\end{result}

\paragraph{Remarks.}
\begin{enumerate}
    \item If the eigenvalues of $X^{\top} X$ are $\kappa_1, \dots, \kappa_p$, then $d_j = \sqrt{\kappa_j}$.
    \item $df(0) = p$, as in least squares.
    \item $df(\lambda) \le df(0)$ for all $\lambda > 0$.
\end{enumerate}

\subsection{Selecting $\lambda$: Generalized Cross-Validation (GCV)}

Define
\[
    \mathrm{GCV}(\lambda)
    = \frac{1}{n} \sum_{i=1}^n
        \left(
            \frac{Y_i - \hat{Y}_i}{1 - \frac{df(\lambda)}{n}}
        \right)^2.
\]

Interpretation: For each $i$, fit the model leaving out observation $i$
and denote the predicted value by $\hat{Y}_{-i}(X_i)$.  
A measure of predictive error is
\[
    \frac{1}{n} \sum_{i=1}^n (Y_i - \hat{Y}_{-i}(X_i))^2.
\]

It can be shown this equals approximately the expression above.

\subsection{Alternate Constrained Form of Ridge Regression}

\[
    \hat{\beta}(\lambda)
    = \arg\min_{\beta}
        \sum_{i=1}^n (Y_i - \beta_0 - \sum_{j=1}^p \beta_j X_{ji})^2
    \quad \text{subject to }
        \sum_{j=1}^p \beta_j^2 \le K.
\]

\subsection{Best Subset Selection}

Choose the ``best'' model with at most $K \le p$ variables:
\[
    \hat{\beta}(K)
    = \arg\min_{\beta}
        \sum_{i=1}^n (Y_i - \beta_0 - \sum_{j=1}^p \beta_j X_{ji})^2
    \quad \text{such that }
        \sum_{j=1}^p 1_{\{\beta_j \ne 0\}} \le K.
\]
This is a \textbf{nonconvex} optimization problem.

\subsection{LASSO (Least Absolute Shrinkage and Selection Operator)}

\[
    \hat{\beta}(K)
    = \arg\min_{\beta}
        \sum_{i=1}^n (Y_i - \beta_0 - \sum_{j=1}^p \beta_j X_{ji})^2
    \quad \text{such that }
        \sum_{j=1}^p |\beta_j| \le K.
\]

Equivalent Lagrangian form:
\[
    \hat{\beta}(\lambda)
    = \arg\min_{\beta}
        \frac{1}{2} \sum_{i=1}^n (Y_i - \beta_0 - \sum_{j=1}^p \beta_j X_{ji})^2
        + \lambda \sum_{j=1}^p |\beta_j|.
\]

\paragraph{Remarks.}
\begin{enumerate}
    \item LASSO uses an $\ell_1$ penalty.  
          The solution is nonlinear in $Y$—no closed form exists.
    \item The problem is convex; efficient solvers exist.
    \item If $K > \sum_j |\hat{\beta}_{LS,j}|$, then LASSO reduces to LS.
    \item The $\ell_1$ penalty induces \textbf{sparsity} in $\hat{\beta}$.
\end{enumerate}

\newpage